\problemname{Counting Islands}
\noindent

Did you know that Sweden is the country with the most islands in the world?
There are about $267{,}000$ in total.
However, many of these islands are very small.

Perhaps it is more relevant to ask which country has the most islands if we only count islands with a land area larger than $100\,\text{m}^2$, or maybe $1000\,\text{m}^2$?

There are $N$ countries, numbered from $1$ to $N$.
Using the internet, you have found the sizes of all islands in these countries.

To satisfy your curiosity, you want to quickly determine which country has the most islands when counting only those with at least $A$ square meters of land area.
You have $Q$ different values of $A$ that you are curious about.


\section*{Input}

The first line of the input contains the integers $N$ and $Q$
($1 \leq N, Q \leq 10^5$),
the number of countries and the number of queries.

The next $N$ lines each describe the islands of one country.
Each of these lines starts with an integer $c$
($1 \leq c \leq 3 \cdot 10^5$),
the number of islands in that country.
This is followed by $c$ integers
$s_1, \dots, s_c$
($1 \leq s_i \leq 10^9$),
the sizes of the country's islands.

Finally, $Q$ lines follow, each containing an integer $A$
($1 \leq A \leq 10^9$), describing the value of $A$ for a query.

Let $T$ be the total number of islands, that is, $T = \sum c$.
It is guaranteed that $T \leq 3 \cdot 10^5$.


\section*{Output}
For each query, print one integer: the index of the country that has the most islands if you only count the ones with a land area of at least $A$
square meters. If two countries are tied for first place (have the same number of islands), then you should answer the one with the higher index.


\section*{Scoring}
Your solution will be tested on a set of test groups, each worth a number of points. Each test group contains
a set of test cases. To get the points for a test group you need to solve all test cases in the test group.

\noindent
\begin{tabular}{| l | l | p{12cm} |}
  \hline
  \textbf{Group} & \textbf{Points} & \textbf{Constraints} \\ \hline
  $1$    & $10$       & $N, T, Q \leq 1000$ \\ \hline
  $2$    & $13$       & $s_i, A \leq 10$ \\ \hline
  $3$    & $16$       & $N \leq 2$ \\ \hline
  $4$    & $22$       & $N \leq 25$ \\ \hline
  $5$    & $39$       & No additional constraints. \\ \hline
\end{tabular}

\section*{Explanation of sample $1$}

In the first query of sample case $1$, country $1$ has three islands with an area of at least $1\,\text{m}^2$: $3, 7, 9$.
Country $2$ has four islands with an area of at least $1\,\text{m}^2$: $1, 2, 2, 3$.
Therefore, country $2$ has the most islands.

In the second query, country $1$ has two islands that are counted, while country $2$ has no islands at all.

\section*{Explanation of sample $2$}

In the first query of sample case $2$, there are no islands at all with an area of at least $4\,\text{m}^2$.
Therefore, all countries have the same number of islands, and the answer is $2$, the country with the highest index.

